\chapter{Testing and Validation}
\label{ch:testing}

\begin{center}
\textit{This chapter presents the testing and validation of the
implementation: the three-tier testing strategy, the unit and integration
suites, security validation, manual end-to-end validation, and the accepted
limitations of testing.}
\end{center}

\vspace*{\fill}
\clearpage
\section{Testing Strategy}

\noindent The project applies a three-tier test strategy:

\begin{enumerate}
  \item \textbf{Unit tests.} Pure-logic modules (guards, policies, validation,
  parsers, bucket planners, renderers, writers) are tested with
  \texttt{node:test} in the API.
  \item \textbf{Database-gated integration suites.} Flows that touch real
  PostgreSQL (academic scope resolution, authorization regression, session
  hardening, OCR worker flow, export scope, job ownership) run as
  \texttt{*.integration.ts} suites against a test database.
  \item \textbf{Manual and end-to-end validation.} A seeded demo dataset drives
  user-observable flows (login, upload, generation, scheduling, attempts,
  exports), recorded in a standing user-validation log.
\end{enumerate}

\noindent Code standards add a fourth gate: \texttt{pnpm typecheck} (strict
TypeScript across all workspaces) and \texttt{pnpm lint} run on the API, web,
and worker sources, and the Python packages are checked with \texttt{ruff}
and \texttt{mypy} (strict mode) plus pytest suites.

\section{API Unit Tests}

\noindent The API suite currently contains 226 tests across 15 suites, all
passing (verified at the time of writing; a full run reports
\texttt{tests 226, suites 15, pass 226, fail 0}). Coverage concentrates on the
non-trivial, security-relevant logic:

\begin{itemize}
  \item \textbf{Authorization and CSRF.} permission catalogue semantics
  (default-deny, \texttt{manage} implication, role/pass composition),
  double-submit CSRF decisions, and cookie attributes;
  \item \textbf{OCR.} coordinator utilities, page inspection/aggregation, and
  worker registry validation;
  \item \begin{sloppypar}\textbf{Material intelligence.} the deterministic
  enhancer (block classification, keep/exclude/review findings, segment
  mapping);\end{sloppypar}
  \item \textbf{Word and worksheet generation.} DOCX and XLSX renderers, PDF
  smoke tests, question/paper document builders;
  \item \textbf{Question planning.} batch bucket building (difficulty
  distribution allocations that sum exactly), batch planning floors/ceilings,
  paper-pattern demand and validation, and the deterministic extractor;
  \item \textbf{Scope resolution and identity.} scope-chain resolution and
  login origin checks.
\end{itemize}

\section{Integration Suites}

\noindent Eleven database-backed integration suites exercise cross-service
behaviour that unit tests cannot:

\begin{itemize}
  \item \texttt{identity/auth-session.integration.ts} --- F1--F5 session
  hardening: session-bound access tokens, strict one-time rotation,
  lineage revocation, logout/revoke semantics, user-status gating, non-
  enumerating login, and password-reset one-shot flow (14 tests);
  \item \texttt{authorization/authz-regression.integration.ts} --- the guard
  chain end to end;
  \item \texttt{\seqsplit{authorization/academic-scope.integration.ts}} and
  \texttt{\seqsplit{resource-scope.integration.ts}} --- scope resolution for teachers and
  students and the 404/403 contract;
  \item \texttt{\seqsplit{authorization/mod-3-export-scope}} /
  \texttt{\seqsplit{mod-4-attempts-scope}}
  and \texttt{\seqsplit{phase-m-remediation.integration.ts}} --- the remediation fixes of
  the Phase~M security audit stay fixed;
  \item \texttt{\seqsplit{authorization/job-ownership.integration.ts}} --- job requester
  stamping and cross-institute adoption defence;
  \item \texttt{academic-structure/*.integration.ts} --- teacher assignments
  and student placements;
  \item \texttt{ocr/ocr-worker-flow.integration.ts} --- the coordinator +
  claim + lease + finalize loop against a real database.
\end{itemize}

\section{Security Validation}

\noindent Security behaviour was validated in a dedicated audit phase that
compared the implementation against a written threat model. It found and
remediated: a HIGH-severity export path that could expose answer keys to
unauthorized teachers (now every export route carries an explicit role
restriction and resolves scope through the authoritative module gate); a
MEDIUM-severity path by which a job created in one institute could be adopted
with another institute's process; and two LOW-severity issues (jobs now record
their requester; stale institute context is cleaned on logout and on
unauthorized responses). Regression suites in the repository keep these fixes
verified.

\section{Manual and End-to-End Validation}

\noindent A seeded demo workload (institutes, academic structure, materials,
syllabi, generated content and questions, patterns, assessments, attempts, and
practice sessions) is exercised through the web client and recorded in the
user-validation log, covering the primary journeys that unit and integration
suites cannot reach: multi-institute sign-in and switching, chunked uploads
through the browser, OCR progress and corrections, AI-generation status
polling, assessment scheduling and timing, attempt submission and the result
view, practice sessions, and printed exports.

\section{Limitations of Testing}

\noindent Browser-level end-to-end automation is not yet present; web
journeys are validated manually against the seeded dataset. OCR accuracy and
AI output quality are validated functionally (structure, schema conformance,
provenance) rather than by quantitative metrics; the AI pipeline's
de-duplication and schema-validation guarantees are covered by worker
unit tests in the Python packages.